\bookchapter{第十二章：影音娱乐与阅读}

\chapterintro{手机现在基本上已经取代了MP3、电视遥控器、Kindle和掌上游戏机。但说实话，很多人只是用了这些App最基本的功能，还有很多好用的地方没发现。这一章我们就来聊聊怎么让手机娱乐更顺手、更舒服。}

\booksection{一、听音乐：选哪个App，怎么听才舒服}

音乐App现在功能都差不多，但各有各的侧重点，选对了能省不少心：

\begin{itemize}
  \item \textbf{网易云音乐}：它的强项是歌单社区和评论区。你想听"适合下雨天听的歌"或者"学累了放松用的纯音乐"，搜一下就有别人整理好的歌单。评论区有时候比歌还精彩。缺点是部分歌曲版权不全，有些歌听着听着就灰了。
  \item \textbf{QQ音乐}：版权最全，周杰伦、五月天这些热门歌手的基本上只有QQ音乐有。音质选项也很丰富，会员可以开无损音质。如果你对"想听的歌都能搜到"这件事很在意，QQ音乐是最稳妥的选择。
  \item \textbf{Apple Music}：苹果用户的首选。界面干净、没有广告、音质好，而且支持无损和杜比全景声。但它的推荐算法不如国内App懂中文用户的喜好，歌单文化也没有网易云那么发达。适合喜欢简洁体验、不太在乎社交功能的人。
  \item \textbf{Spotify}：全球用户最多，推荐算法公认最强。"每周发现"和"日推"准确率高得吓人，用越久越懂你。缺点是免费版广告多，而且某些中文歌版权覆盖不全。需要特殊网络才能用。
\end{itemize}

\textbf{几个听歌小技巧}：如果你用的是蓝牙耳机，去音乐App设置里看看有没有“音效调节”或“均衡器”，稍微调一下低音和高音，效果提升明显。另外，流媒体听歌很费流量，记得在设置里把“仅WiFi下播放无损”打开，或者提前下载好喜欢的歌单离线听。

\booksection{二、播客（Podcast）：通勤路上的“耳朵伴侣”}

播客这个东西，没听过的人觉得跟自己无关，一旦开始听就回不去了。本质就是“电台节目的Netflix版”——有人聊科技，有人讲历史，有人读故事，你想听什么类型的都有。

\begin{itemize}
  \item \textbf{小宇宙App}：国内做得最好的播客App，界面清爽，没有乱七八糟的功能。它的"每日推荐"和"编辑精选"帮你发现新节目，评论区也很活跃。我个人听的最多的是《忽左忽右》（文化历史类）和《不合时宜》（社会议题类），每期都像在旁听几个聪明人聊天。
  \item \textbf{苹果播客}：iPhone自带，不需要额外下载App。和CarPlay的集成最好，开车的时候很方便。
  \item \textbf{收听技巧}：大部分播客App都支持1.2倍速或1.5倍速播放，习惯了之后能省不少时间。我一般1.2倍听中文，1.5倍听英文。另外一定要开"自动下载新节目"和"播完自动删除"，这样你的手机不会堆满听过的音频文件。
\end{itemize}

\booksection{三、看视频：从播放器到流媒体的正确打开方式}

很多人不知道，手机上可以播的格式远比流媒体App里提供的多。一个好的本地播放器能解决很多问题。

\begin{itemize}
  \item \textbf{VLC播放器}：开源免费，几乎什么格式都能播——MKV、AVI、FLAC、甚至是ISO光盘镜像。如果你下载了一个视频在电脑上能播但手机上播不了，八成是格式问题，装个VLC就解决了。它还支持外挂字幕，把.srt字幕文件和视频放在同一个文件夹，VLC会自动加载。
  \item \textbf{nPlayer}：如果你想要更精致的体验，这个付费App值得买。界面好看，手势控制很顺手——左边上下滑调亮度，右边上下滑调音量，左右滑快进快退。支持DTS和杜比音效，看大片的时候临场感强很多。
  \item \textbf{Bilibili（B站）}：已经不只是二次元网站了，现在上面的学习资源多到离谱。考研、编程、摄影、烹饪，什么教程都有。弹幕文化是B站的灵魂，看视频的时候开着弹幕，感觉像有一群人在陪你一起看。不过看学习视频的时候建议关弹幕，不然容易分心。
  \item \textbf{投屏到电视}：手机屏幕再大也不如电视看着舒服。同一个WiFi下，打开视频App右上角的"TV"投屏按钮，就能把画面投射到电视上。大部分智能电视和投影仪都支持，看剧看电影体验好太多。iPhone用户还可以用AirPlay直接镜像整个屏幕。
\end{itemize}

\booksection{四、电子书阅读：手机上也能好好看书}

有人说手机上读书伤眼睛、容易分心，这些确实有道理。但手机阅读的优势也很明显——随时随地，碎片时间利用最大化。我在地铁上、排队的时候、睡前，用手机读了不少书。

\begin{itemize}
  \item \textbf{微信读书}：目前国内体验最好的阅读App。书库很全，出版书和网文都有。它用"阅读时长换书币"的模式，看得多就能免费换书，对学生党很友好。它的"划线"功能可以看别人在哪些段落做了标记，有时候看到别人标注的好句子，确实有共鸣。不过"阅读排行榜"这个东西有点卷，容易变成"为了刷时长而读书"，建议关掉。
  \item \textbf{Kindle App}：如果你有Kindle设备，手机端Kindle App可以同步阅读进度。Kindle的正版电子书排版质量比盗版PDF好太多，但Kindle中国区书城已经关了，需要切换到海外区。
  \item \textbf{Apple Books（图书）}：iPhone自带，适合读EPUB和PDF。界面极简，翻页动画很舒服。如果你从网上下载了电子书文件，直接"分享到图书"就能导入。
  \item \textbf{阅读设置建议}：不管用哪个App，一定要把背景设成暖黄色调（护眼模式），字体调大一点，行间距拉开。晚上读书打开"深色模式"——黑底白字，眼睛舒服很多。千万别白底黑字小字体在暗光下读太久，眼睛疲劳是积累的。
\end{itemize}

\booksection{五、手机游戏：玩得爽，也不伤手机}

我对手机游戏的态度是：玩可以，但别让它反过来控制你。合理使用游戏工具，玩得更舒服，也不耽误正事。

\begin{itemize}
  \item \textbf{游戏模式/游戏空间}：现在主流手机都有这个功能。打开之后手机会自动清后台、屏蔽通知、防止误触——你正在打团战，突然一个电话弹出来，这种事情游戏模式能帮你挡掉。部分手机还支持性能模式/E-sports模式，强制处理器跑满，但代价是发热和耗电都更快。
  \item \textbf{画质和帧率怎么调}：不要无脑拉满。先保证帧率（流畅度），再追求画质。60帧比30帧舒服太多了，但120帧对手机压力很大，中端机不用强求。如果你玩的是竞技类游戏（比如王者、吃鸡），帧率优先级远高于画质——流畅的操作手感比好看的草皮重要一百倍。
  \item \textbf{游戏时手机发烫怎么办}：这是正常的，但长期高温会损伤电池寿命。几个降温办法：摘掉手机壳散热；把屏幕亮度降到70\%左右（亮度是高发热源）；不要一边充电一边打游戏（发热叠加，电池死得快）。如果手机烫到手握不住的程度，赶紧退出来让它歇一会儿。
  \item \textbf{防沉迷小工具}：iOS的"屏幕使用时间"和安卓的"数字健康"都可以给游戏App设每日时间限额。我给自己设的是工作日最多一小时，周末两小时。到时间了App图标会变灰，点不进去。这个限制是真的有用，不是装装样子的。
  \item \textbf{云游戏}：如果你手机性能不够但又想玩大作，可以试一下云游戏平台（比如腾讯START云游戏、网易云游戏）。游戏在远程服务器上跑，画面以视频流的方式传到你的手机上。对网络要求高但对你手机性能几乎没有要求。缺点是延迟感不可避免，适合玩策略类、RPG类，不适合FPS射击类。
\end{itemize}

\begin{tcbpractice}{来动手：打造你的手机娱乐系统}
\begin{enumerate}
  \item 在音乐App里下载你最喜欢的3张专辑或歌单，确保离线可听。
  \item 去小宇宙或苹果播客订阅3个感兴趣的频道，下载一期试听。
  \item 安装VLC播放器，尝试用它播放一段电脑拷贝过来的视频文件。
  \item 在微信读书里找一本想读的书，设置暖色背景和护眼模式，把字体调到适合自己的大小。
  \item 如果你玩游戏，打开游戏模式看看有哪些设置可以优化（免打扰、帧率优先等）。
\end{enumerate}

\end{tcbpractice}
